\subsection{The full stable norm condition at an allocated prime}

\begin{statement}{A.131}{Proved}
{The exact tame norm tower and its retained branch character}
Let \(q\ne3\) be prime, let \(F=\Q_q\), and put
\(s=v_3(q-1)\).  Choose a generator \(g\) of \(\F_q^\times\)
and let
\[
 \lambda_s:\mathcal O_F^\times\longrightarrow\Z/3^s\Z
 \tag{A.131.1}
\]
be reduction followed by the character sending \(g\) to one.
For \(s=0\) its target is the zero group.
Fix \(z=q^h u\), with \(u\in\mathcal O_F^\times\) and
\(h\in\Z\) prime to three.  Choose compatible roots and set
\[
                F_j=F(\theta_j),\quad
              \theta_j^{3^j}=z,\quad \theta_{j+1}^3=\theta_j.
 \tag{A.131.2}
\]
For \(b=q^t v\in F^\times\), put
\(c=\lambda_s(v)\) and \(a=\lambda_s(u)\).
The exact norm criterion at every level is
\[
 b\in N_{F_j/F}F_j^\times
       \quad\Longleftrightarrow\quad
                 hc-ta=0\pmod{3^{\min(j,s)}}.
 \tag{A.131.3}
\]
The norm images therefore stabilize at level \(s\), with
the full all-level condition \(hc-ta=0\bmod3^s\).
Moreover the image of the actual compatible norm family is
\[
 \operatorname{im}\left(\varprojlim_j F_j^\times
                      \xrightarrow{\ N\ }F^\times\right)
       =\{q^t v: h\lambda_s(v)-t\lambda_s(u)=0\bmod3^s\}.
 \tag{A.131.4}
\]
The inverse system uses the actual field norms.
The same assertion holds for the pro-three completions,
with \(t\in\Z_3\) and the literal unit coordinate in
\(\Z/3^s\).

Here is its map to the norm-chain update.
Suppose \(N=3^n\), \(n\ge1\), \(q\) splits
\(\Q(a^{1/N},\mu_N)\), and \(q\nmid a\).
Label the actual roots in \(F\) by
\(\alpha_i=\zeta_N^i\alpha_0\), \(0\le i<N\).
Let \(Q=q^kQ_0\), \(Q_0,R\in\mathcal O_F^\times\),
and \(e,k\in\Z\) with \(3\nmid ek\).
For local target coordinates \(\beta_i\), put
\[
 \begin{gathered}
 t_i=v_q(\beta_i),\quad
 c_i=\lambda_s(q^{-t_i}\beta_i),\quad
 a_i=\lambda_s(\alpha_i),\\
 x=\lambda_s(Q_0),\qquad y=\lambda_s(R).
 \end{gathered}
 \tag{A.131.5}
\]
The target \(\beta_iR^e\) has an actual compatible family
of norm preimages in all extensions
\(F((Q^e\alpha_i)^{1/3^j})\) for every \(i\) exactly when
\[
       ek(c_i+ey)-t_i(ex+a_i)=0\pmod{3^s}
                              \quad(0\le i<N).
 \tag{A.131.6}
\]
Equivalently, division by the units \(e,k\) gives the
literal coordinate equation
\[
             c_i=\frac{x}{k}t_i+\frac{a_i}{ek}t_i-ey
                                        \quad\text{in }\Z/3^s.
 \tag{A.131.7}
\]
The branch character in this formula satisfies
\[
       a_i=a_0+i\kappa3^{s-n}\pmod{3^s},
                   \qquad 3\nmid\kappa.
 \tag{A.131.8}
\]
In particular the \(N\) roots cannot all have zero unit
character at the full stable precision \(3^s\).
Splitting a deeper root field may make their characters
zero modulo a chosen \(3^r\); the actual branch difference
then requires \(r\le s-n\).

For a relative norm-one local target, write \(d=N/3\)
and suppose its coordinates have product one in each
block \(h,h+d,h+2d\).  The full equations imply
\[
   \kappa3^{s-1}
          (t_{h+d}+2t_{h+2d})=3e^2ky\pmod{3^s}
                                \quad(0\le h<d).
 \tag{A.131.9}
\]
For \(R=1\) this forces all three valuations in each
block to be equal modulo three.  With arbitrary scalar
\(R\), their displayed block moments must agree modulo
three; when \(s=1\) they must all be zero.
These are consequences of the exact norm equations for
the stated norm-one target.  For a general existing
lower representative its unit and valuation coordinates
must instead be retained in the block sum of (A.131.6).

There is a further exact comparison with cyclic repair.
For the same relative norm-one target, suppose its
coordinates at the full stable precision satisfy
\(c_i=z_0t_i+w\) with common scalars \(z_0,w\in\Z/3^s\).
If (A.131.6) also holds, then
\[
                           t_i=0\pmod3\quad\text{for every }i.
 \tag{A.131.14}
\]
Thus a primitive inertia vector in one of the scalar-shifted
cyclic graphs cannot simultaneously satisfy this full
branch norm condition for the stated norm-one target.
The actual graph needed by the norm equation remains
(A.131.7), including its diagonal branch term.
\end{statement}

\begin{proof}
\emph{The fields, uniformizers and complete norm image.}
Since \(h\) is prime to three, the valuation of
\(\theta_j\) over \(F\) has denominator \(3^j\).
It follows that \([F_j:F]=3^j\), and the extension is
totally ramified, of residue degree one.
It is tame since \(q\ne3\).
Choose integers \(A_j,B_j\) satisfying
\(A_jh+B_j3^j=1\).
The actual element
\[
       \pi_j=\theta_j^{A_j}q^{B_j}
               \quad\text{has normalized }F_j\text{-valuation one},
       \qquad N_{F_j/F}\pi_j=q u^{A_j}.
 \tag{A.131.10}
\]
The second identity uses
\(N\theta_j=z\), whose sign is positive because
the degree \(3^j\) is odd, and \(Nq=q^{3^j}\).
It is valid also for negative \(h,A_j,B_j\).

For every unit of \(F_j\), its residue norm is its
\(3^j\)th power in the common residue field \(\F_q\).
Conversely every such residue norm is supplied by a
unit of \(F\) itself.  The principal-unit norm is
onto \(1+q\Z_q\): the power map
\(w\mapsto w^{3^j}\) is an automorphism of that
group, by Hensel lifting, and is exactly the norm
of a scalar \(w\in F\subset F_j\).
It follows that the unit norm group is precisely
\[
 N_{F_j/F}\mathcal O_{F_j}^\times
     =\{v\in\mathcal O_F^\times:
                         \overline v\in(\F_q^\times)^{3^j}\}.
 \tag{A.131.11}
\]
Every element of \(F_j^\times\) is \(\pi_j^t\)
times a unit, and its norm has valuation \(t\),
because the residue degree is one.
Equations (A.131.10)--(A.131.11) show that
\(q^t v\) is a norm exactly when
\(v u^{-A_jt}\) has \(3^j\)-power residue.
In the cyclic residue group this means
\(c-A_jta=0\bmod3^{\min(j,s)}\).
The identity \(A_jh\equiv1\bmod3^j\) proves
(A.131.3).  This direct norm computation also
agrees with the general local norm limitation
theorem \cite[III, Theorem 3.5, pp.~108--109]{MilneCFT};
it does not replace a non-Galois field by its closure.

\emph{Actual compatibility of the local preimages.}
Necessity in (A.131.4) follows by projection to every
level.  Conversely suppose (A.131.3) holds at all
levels, and for each \(j\) consider
\[
       X_j=\{w\in F_j^\times:N_{F_j/F}w=b\}.
 \tag{A.131.12}
\]
This set is nonempty by the proved criterion.
Every member has normalized valuation \(t=v_q(b)\).
Consequently it is a closed subset of the compact
coset \(\pi_j^t\mathcal O_{F_j}^\times\), hence compact.
The actual transition norm maps send \(X_{j+1}\)
into \(X_j\) by transitivity, and are continuous.
Every finite initial set of compatibility equations
has a solution: choose a member of its largest
\(X_J\) and take its field norms to all lower levels.
The product of the compact \(X_j\), with these
closed equations, therefore has the finite intersection
property and a nonempty intersection.
This gives an actual compatible family, not merely
separate preimages at individual levels.

For the completed assertion, the pro-three
completion of a local multiplicative group of
residue characteristic \(q\ne3\) is the product
of \(\Z_3\) and the three-primary residue roots
of unity.  The norm formulas above extend through
these actual quotient maps; the principal-unit
factor has trivial pro-three completion.
Thus its image is still exactly (A.131.3), now
with \(t\in\Z_3\).
The completed groups themselves are compact.
Their nonempty norm fibres and the same finite
intersection argument prove the completed version
of (A.131.4).

\emph{The update and the branch character.}
The radicand \(Q^e\alpha_i\) has valuation \(ek\)
and unit part \(Q_0^e\alpha_i\), whose character is
\(ex+a_i\).
The target \(\beta_iR^e\) has valuation \(t_i\)
and unit character \(c_i+ey\).
Substituting these in the proved full norm criterion
gives (A.131.6) and its equivalent (A.131.7).
Since only finitely many branches occur, the
compatible family argument can be applied to every
one and then their finite product.

The chosen primitive root \(\zeta_N\) reduces
to a residue root of exact order \(3^n\).
Its image under \(\lambda_s\) therefore has
the exact form \(\kappa3^{s-n}\), with \(\kappa\)
prime to three.  The equality
\(\alpha_i=\zeta_N^i\alpha_0\) proves (A.131.8).
This difference is nonzero at full precision.
For it to be zero modulo \(3^r\), one needs
\(r\le s-n\), proving the asserted restriction
on an attempted removal of the branch term.

\emph{The relative block consequence.}
For the stated norm-one target the block product
equalities give
\(\sum_{i\equiv h\bmod d}t_i=0\) exactly and
\(\sum_{i\equiv h\bmod d}c_i=0\bmod3^s\).
Summing (A.131.6) in that block leaves
\[
             3e^2ky=\sum_{i\equiv h\bmod d}a_it_i.
 \tag{A.131.13}
\]
The constant part of \(a_i\), including the term
\(h\kappa3^{s-n}\), cancels by the zero sum of
valuations.  The remaining value is
\(\kappa3^{s-n}d(t_{h+d}+2t_{h+2d})\), which is
(A.131.9) because \(d=3^{n-1}\).
For \(y=0\), its vanishing is precisely
\(t_{h+d}+2t_{h+2d}=0\bmod3\).
Together with the zero valuation sum it says
that all three valuations are equal modulo three.
For arbitrary \(y\) the right side of (A.131.9)
is the same for every block, so the block moments
agree; at \(s=1\) it is zero, giving the additional
assertion.  This derives every stated moment
condition through the actual norm and unit maps.

\emph{The primitive-inertia comparison.}
Substitute \(c_i=z_0t_i+w\) in (A.131.6).
Within one block, renumber its three coordinates by
\(i=0,1,2\) and put \(\delta=\kappa3^{s-1}\).
The resulting equations have precisely the form
\[
                  (u+i\delta)t_i=C\pmod{3^s},
                   \qquad t_0+t_1+t_2=0\pmod{3^s},
 \tag{A.131.15}
\]
where \(u=a_h-e(kz_0-x)\) and
\(C=ek(w+ey)\).
Suppose first that \(v=v_3(u)<s-1\).
Each equation forces \(3^v\mid C\).
Divide by \(3^v\), writing \(u'=u/3^v\),
\(C'=C/3^v\) and \(\delta'=\delta/3^v\), and work
modulo \(3^{s-v}\).
Here \(u'\) is a unit, \(s-v\ge2\), and
\((\delta')^2=0\) in the coefficient ring.
Consequently
\[
       t_i=C'\bigl((u')^{-1}-i\delta'(u')^{-2}\bigr),
                \qquad \sum_{i=0}^2 t_i=3C'(u')^{-1}
                         \pmod{3^{s-v}}.
 \tag{A.131.16}
\]
The second identity uses \(3\delta'=0\).
The zero sum gives \(3^{s-v-1}\mid C'\), so every
\(t_i\) is divisible by three.

If instead \(3^{s-1}\mid u\), every equation first
forces \(3^{s-1}\mid C\).
Divide by \(3^{s-1}\) and reduce modulo three.
The three coefficients are \(u'+i\kappa\).
Exactly one is zero; its equation forces \(C'=0\).
The other two valuations are zero modulo three,
and their zero sum forces the third to be zero.
This includes \(s=1\).
Every block therefore gives (A.131.14).
All the scalars and equations used here come from
the displayed local norm map; no restriction is
deduced for an unspecified lower representative.
\end{proof}

The retained exact program
\path{code/stable_local_norm_certificate.py}
checks 7,980 literal residue-norm cosets, 31,920
uniformizer norm-transition identities, 1,890 stable
intersections, 4,401 relative block moments and 110
branch-depth comparisons.  It also checks the primitive
inertia comparison in 551,880 literal systems through
stable precision \(3^4\).  Its saved output includes
the primes and precisions used; compact compatibility
is established by the proof, not by these finite checks.
